\documentclass[a4paper]{article}
\usepackage{amsfonts}
\usepackage{amsmath}
\usepackage{amssymb}
\usepackage{graphicx}     

\begin{document}
  
\noindent \large Auteur: Stan Van Campenhout \\
\noindent \large Studierichting:  Informatica\\
\noindent \large Oefeningenreeks 4A nr. 23.2\\

\medskip

\normalsize

$ \displaystyle\int x \operatorname{Bgtan} x dx = \dfrac{1}{2}\displaystyle\int\operatorname{Bgtan}x d(x^2) $\\

$ \dfrac{x^2}{2}\operatorname{Bgtan}x - \dfrac{1}{2}\displaystyle\int x^2 d\left(\operatorname{Bgtan}x\right)
= \dfrac{x^2}{2}\operatorname{Bgtan}x - \dfrac{1}{2}\displaystyle\int \dfrac{x^2}{1 + x^2} dx $\\

$ = \dfrac{x^2}{2}\operatorname{Bgtan}x - \dfrac{1}{2}\displaystyle\int \dfrac{1 + x^2 - 1}{1 + x^2} dx $\\

$ = \dfrac{x^2}{2}\operatorname{Bgtan}x - \dfrac{1}{2}\displaystyle\int \dfrac{1 + x^2}{1 + x^2} dx 
+ \dfrac{1}{2}\displaystyle\int \dfrac{1}{1 + x^2} dx $\\

$ = \dfrac{x^2}{2}\operatorname{Bgtan}x - \dfrac{1}{2}x + \dfrac{1}{2}\operatorname{Bgtan}x + c  
 = \dfrac{x^2 + 1}{2}\operatorname{Bgtan}x - \dfrac{1}{2}x + c $\\

\begin{thebibliography}{999}
%\bibitem{a} Referentie
\end{thebibliography}
\end{document} 
